\section{Introduction}
\label{sec:intro}

Reinforcement learning (RL) has emerged as a cornerstone of autonomous decision-making in sequential control tasks and interactive video games~\cite{mnih2015human,sutton2018reinforcement}. Classic 2D platformers represent a formidable testbed because they demand continuous temporal precision, rapid reaction to moving and stationary hazards, and robust spatial navigation. Unlike discrete grid-worlds, platformers feature continuous kinematic properties, requiring jump timing calibrated to gravitational acceleration, obstacle heights, and pit widths.

In this university team project, we design, implement, and evaluate an autonomous game agent navigating a 2D graveyard platformer entitled \textit{Ho{\'a}'s Midnight Escape}. In this scenario, the protagonist Ho{\'a} awakens in an eerie midnight graveyard and must navigate through tall tombstones, open burial pits, and airborne hazards to reach the safety of his home. Initially, the agent possesses no prior domain knowledge; through trial-and-error interaction and environmental feedback, it acquires an optimal control policy.

The primary motivation of this work is to demonstrate that an interpretable, resource-efficient Tabular Q-Learning agent can achieve high navigational proficiency in continuous 2D games when coupled with domain-specific spatial discretization and structured curriculum learning. To prevent merge friction and maintain software engineering rigor across our four-member team, our application decouples game physics, state representation, reward formulation, and empirical evaluation via explicit abstract interfaces.

The remainder of this paper is structured as follows: Section~\ref{sec:related} reviews related literature in game AI and reinforcement learning. Section~\ref{sec:objectives} outlines project objectives and technical gaps. Section~\ref{sec:teamwork} describes team task distribution, software workflow, and project management. Section~\ref{sec:approach} details the input/output representation, Q-learning algorithms, and modular reward strategies. Section~\ref{sec:impl} discusses system architecture, UI design, tech stack, and code highlights. Section~\ref{sec:eval} presents an extensive empirical evaluation with statistical significance tests across curriculum tiers. Finally, Section~\ref{sec:conclusion} concludes the paper with reflections and future research directions.
